% ============================================================
% Chapter 4 — Benchmark Construction and Human Annotation Protocol
% Draft for master thesis body
% ============================================================
%
% Usage: paste into thesis main .tex as \chapter{Benchmark Construction and Human Annotation Protocol}
% Dependencies: booktabs, multirow, array, longtable, graphicx, listings (for JSON snippet)
% Estimated length: 12–15 pages
% ============================================================

\chapter{Benchmark Construction and Human Annotation Protocol}
\label{chap:benchmark}

This chapter specifies how the evaluation material used throughout the
thesis was assembled and how human reference labels were obtained.
Section~\ref{sec:bench-goals} states the design goals that the benchmark
must satisfy. Sections~\ref{sec:track-a-construction}
and~\ref{sec:track-b-construction} document the construction of the two
tracks introduced in Chapter~\ref{chap:methodology}: the screenshot-based
reproduction track that reuses UIClip resources, and the
requirement-driven generation track built primarily on a reduced
Vision2Web Level~1/Level~2 subset, with Design2Code retained as a
static UI-to-code reference point. Section~\ref{sec:item-schema} defines the
on-disk item schema and the result-submission format consumed by the
public leaderboard. Section~\ref{sec:rubric} introduces the four-dimensional
GUI quality rubric, including the anchor descriptions that operationalise
each scoring level. Sections~\ref{sec:annotation-protocol}
and~\ref{sec:irr} describe the human annotation procedure and the
inter-rater reliability targets used to validate it. Finally,
Section~\ref{sec:pilot} reports the pilot annotation that calibrated the
rubric before the main collection.

% ----------------------------------------------------------------
\section{Benchmark Design Goals}
\label{sec:bench-goals}
% ----------------------------------------------------------------

The benchmark is designed around four goals.

\paragraph{Reproducibility.} Every item, prompt, model identifier, and
random seed used in any reported experiment is recorded so that the same
numbers can be regenerated by a third party. Each Track~A item retains
its provenance back to the original UIClip release; each Track~B item
retains the requirement text, generator model identifier, and rendering
parameters used to produce it.

\paragraph{Extensibility.} New models must be evaluable without
modifying the benchmark. The on-disk item schema
(Section~\ref{sec:item-schema}) is fixed, the rubric is fixed
(Section~\ref{sec:rubric}), and the result-submission format separates
\emph{judge identity} (provider, model, prompt strategy) from
\emph{benchmark content}. Adding a new judge therefore requires only
running the public evaluation script and submitting a JSON file.

\paragraph{Track separation with shared semantics.} Track~A
(screenshot-only) and Track~B (executable) keep distinct directory
trees, item schemas, and reference label types but share the rubric and
the result-submission format. This allows cross-track comparison while
preventing accidental conflation of derived UIClip captions with
human-rated rubric scores.

\paragraph{Leaderboard compatibility.} A minimal Hugging Face Space
dashboard reads the result JSON files directly from the project
repository and ranks generated GUI submissions, aggregated by generator
model and requirement set. New entries reach the leaderboard through
pull requests rather than online inference, keeping the operational
surface small and the cost of hosting near zero. Judge models and human
labels are used to compute metrics, but the leaderboard's primary
object is the quality of the generated interfaces. The implementation
is documented in Appendix~\ref{appendix:leaderboard}.

% ----------------------------------------------------------------
\section{Track A: UIClip Reproduction Subset}
\label{sec:track-a-construction}
% ----------------------------------------------------------------

\subsection{Source}

Track~A reuses the public UIClip human caption dataset
(\texttt{biglab/uiclip\_human\_data\_hf} on Hugging Face), which contains
single-screenshot rows tagged with natural-language captions of the form
\textit{``UI screenshot. well-designed. \{page description\}''} or
\textit{``UI screenshot. poor design. bad contrast. \{page description\}''}.
No additional collection is performed; the dataset is treated as a
fixed external resource.

\subsection{Caption-Matched Pair Construction}

The dataset does not provide explicit \texttt{img\_good} /
\texttt{img\_bad} pair columns. Pairs are derived deterministically by:
(i)~classifying each row as \emph{good} or \emph{bad} based on the
presence of \textit{well-designed} versus \textit{poor design / bad
contrast / bad proximity} in the caption;
(ii)~normalising the page description (lowercasing, removing the
quality keywords and articles, collapsing whitespace);
(iii)~matching one good row with one bad row that share the same
normalised description. This yields what are referred to throughout the
thesis as \emph{caption-matched derived pairs}, distinct from raw
human-judged pairwise preference labels.

This construction caveat is stated explicitly in every results table
that uses Track~A so that readers do not interpret the labels as
direct human preference judgements.

\subsection{Sampling Plan and Manual Cleaning}

Fifty pairs are drawn at \texttt{seed=42} from the first 2{,}500 source
rows; if scaled later, the same seed and source-row cap are reused so
that the smaller sample is a strict prefix of the larger one. Each pair
is briefly inspected by the author to flag mismatches in which the
\emph{good} and \emph{bad} screenshots clearly do not depict the same
underlying page despite sharing a description. In the pilot batch this
removed seven pairs out of fifty (item ids 0, 15, 31, 33, 35, 41, 48),
yielding a clean subset of 43~pairs that is used as the canonical
Track~A sample for all multi-condition comparisons in
Chapter~\ref{chap:experiments}.

\subsection{Item Storage}

\begin{lstlisting}[basicstyle=\ttfamily\small]
track_a/
  pair_0001/
    img_a.png        % the "good" screenshot
    img_b.png        % the "bad"  screenshot
    label.json       % {"truth_position": "A", "source": "uiclip_human",
                     %  "group_key": "...", "captions": {...}}
\end{lstlisting}

No new annotations are collected for Track~A. The reference label is
the \texttt{truth\_position} field, derived from the
well-designed/poor-design caption distinction.

% ----------------------------------------------------------------
\section{Track B: Requirement-Driven Generated Interfaces}
\label{sec:track-b-construction}
% ----------------------------------------------------------------

Track~B supplies the executable interfaces required by the rubric and
dynamic evaluations. Each item is a self-contained directory containing
a natural-language requirement, the HTML produced by a generator model,
a rendered screenshot, and a serialised DOM snapshot.

\subsection{Requirement Sources}

\paragraph{Source 1: Vision2Web Level~1/Level~2 curation.}
Vision2Web~\cite{vision2web} is used as the primary Track~B source
because it already frames website generation as a hierarchy from static
visual implementation to interactive frontend behaviour and full-stack
development. This structure aligns closely with the thesis objective of
comparing static GUI-quality signals with dynamic task-validation
signals on the same generated interfaces.

The thesis uses a reduced subset of Level~1 static webpage tasks and
Level~2 interactive frontend tasks. Level~1 items provide controlled
visual prototypes and requirements for static rubric scoring. Level~2
items additionally provide interaction-oriented requirements and
testable functional expectations, making them suitable for the dynamic
validation procedure in Chapter~\ref{chap:dynamic}. Level~3 full-stack
tasks are excluded because backend state, deployment, authentication,
and long-horizon workflow requirements would shift the project from GUI
quality evaluation into full-stack software engineering evaluation.

The pilot subset contains approximately 10--20 tasks. Items are selected
for clear visual requirements, local renderability, and the presence of
at least one short interaction task with an unambiguous success
criterion. If the pilot is successful, the formal subset may be expanded
within the time and annotation budget.

\paragraph{Source 2: Design2Code and author-written reserve items.}
Design2Code~\cite{design2code} is retained as a static UI-to-code
reference point and as a reserve source for simple reproduction cases.
It is no longer the main Track~B source because it is primarily a
screenshot-to-code benchmark and does not provide the same direct link
between visual quality and functional validation as Vision2Web. A small
set of author-written form-heavy requirements may be added only if the
selected Vision2Web subset lacks enough clear dynamic-validation cases.
These reserve requirements are kept intentionally simple so that
failures can be attributed to interface quality rather than ambiguous
task wording.

\subsection{Generation}

Each requirement is given to two or three generator models drawn from
the same lineup as the judges (Section~\ref{sec:model-lineup} in
Chapter~\ref{chap:experiments}). Generators are prompted to return a
single self-contained HTML file with inline CSS and minimal JavaScript;
external network requests, build steps, and framework dependencies are
disallowed so that every output can be rendered directly. The generator
identity is recorded in the item metadata and is not revealed to either
judges or human annotators during evaluation.

\subsection{Rendering Pipeline}

Each generated HTML file is rendered with headless Chromium via
Playwright at a fixed viewport (1280$\times$800 for desktop pages,
390$\times$844 for mobile). The pipeline produces three artefacts: a
full-page PNG screenshot used by static evaluators; a serialised DOM
snapshot (a recursive JSON dump of element tag, accessible name,
relevant attributes, and bounding box) used by the dynamic evaluator
in Chapter~\ref{chap:dynamic}; and a render log capturing console
errors and unresolved resources, used to flag items that fail to load.

\subsection{Item Storage}

\begin{lstlisting}[basicstyle=\ttfamily\small]
track_b/
  item_0001/
    requirement.md       % the natural-language UI requirement
    generated.html       % the LLM-generated implementation
    screenshot.png       % full-page render at fixed viewport
    dom.json             % serialised element tree
    generator_meta.json  % {"model": "...", "prompt_id": "...",
                         %  "viewport": [W, H], "render_log": "..."}
\end{lstlisting}

% ----------------------------------------------------------------
\section{Benchmark Item Schema and Result Format}
\label{sec:item-schema}
% ----------------------------------------------------------------

The benchmark commits to two on-disk schemas: an \emph{item schema} that
describes a single evaluation unit, and a \emph{result schema} that a
judge submits for the entire benchmark. Both are documented in full in
Appendix~\ref{appendix:schema}; only the essential structure is shown
here.

\subsection{Item Schema}

Every item, regardless of track, has the following common fields
augmented by track-specific blocks:

\begin{lstlisting}[basicstyle=\ttfamily\small]
{
  "item_id":   "track_b/item_0017",
  "track":    "B",
  "rubric_dimensions": ["D1", "D2", "D3", "D4"],
  "human_reference": {
    "rubric_scores": {"D1": 4, "D2": 3, "D3": 5, "D4": 2},
    "n_annotators": 3,
    "irr_alpha":  {"D1": 0.62, "D2": 0.58, "D3": 0.71, "D4": 0.49}
  },
  "track_b": {
    "requirement_md":  "...",
    "generator_model": "gpt-4o-2024-08-06",
    "viewport":        [1280, 800]
  }
}
\end{lstlisting}

Track~A items omit the \texttt{track\_b} block and replace
\texttt{human\_reference} with the derived \texttt{truth\_position}.

\subsection{Result Submission Format}

A submission is one JSON file containing a metadata block (evaluator
provider, model, prompt strategy, generator model, run date, repository
commit) and a flat list of per-item results. For Track~A each result
contains a pairwise choice used for baseline analysis. For Track~B each
result contains static rubric scores, dynamic validation outcomes, and
the generated interface metadata needed to aggregate scores by
generator model. The leaderboard renders these submissions directly
without re-executing the evaluator or generator. A minimal example and the
JSON Schema are provided in Appendix~\ref{appendix:schema}.

% ----------------------------------------------------------------
\section{GUI Quality Rubric}
\label{sec:rubric}
% ----------------------------------------------------------------

\subsection{Design Rationale}

The rubric translates established usability theory into a small number
of dimensions that human annotators and LLM-based evaluators can score
reliably from a screenshot (and, for Track~B, an executable page).
Two design constraints shape it. First, the dimensions must be
\emph{distinguishable}: two annotators reading the anchors should agree
on which dimension a given criticism belongs to. Second, the dimension
count must be small enough that one annotator can score all of them on
one item in three to five minutes; pilot timing data
(Section~\ref{sec:pilot}) confirms this budget.

The conceptual roots are Nielsen's heuristics for user interface
design~\cite{nielsen1990heuristic, nielsen1994usability} and the ISO
9241-11 definition of usability~\cite{iso9241-11}. Rather than
replicating Nielsen's ten heuristics directly, four of which (system
status visibility, user control, help and documentation, error
prevention) are difficult to assess from a static screenshot, the
rubric folds the screenshot-observable subset into two static
dimensions and adds two Track-B-specific dimensions for
requirement satisfaction and interactive behaviour.

\subsection{The Four Dimensions}

\begin{description}
  \item[D1 — Visual Structure.] Layout, alignment, spacing, and the
    visual hierarchy that guides the eye through the interface.
    Captures the screenshot-observable subset of Nielsen's
    \emph{Aesthetic and Minimalist Design} and \emph{Consistency and
    Standards}.

  \item[D2 — Information Clarity.] Grouping of related items, label
    quality, readability of text, and the ease of inferring the purpose
    of each region without prior context. Reflects ISO 9241-11
    \emph{learnability} from the visual surface alone.

  \item[D3 — Requirement Fidelity \textit{(Track~B only)}.] The degree
    to which the rendered interface implements the elements,
    structure, and functionality stated in its natural-language
    requirement. The reference for this dimension is the original
    requirement text, not annotator preference.

  \item[D4 — Interaction Quality \textit{(Track~B only, dynamic)}.]
    Responsiveness to user actions, appropriateness of feedback and
    state transitions, and the smoothness of completing a derived
    task. Scored only after running the dynamic-evaluation procedure
    of Chapter~\ref{chap:dynamic}; not assessable from a static
    screenshot alone.
\end{description}

All dimensions use the same 1--5 ordinal scale.

\subsection{Anchor Descriptions}

Each anchor is a short prose description grounded in observable
features. The complete anchor table is reproduced verbatim in the
annotator guide (Appendix~\ref{appendix:annotation}); the abbreviated
form below shows the conceptual span of each dimension.

\begin{table}[ht]
\centering
\small
\begin{tabular}{@{}p{2.6cm} p{3.6cm} p{3.6cm} p{3.6cm}@{}}
\toprule
\textbf{Dimension} & \textbf{1 — Poor} & \textbf{3 — Acceptable} & \textbf{5 — Excellent} \\
\midrule
\textbf{D1} Visual Structure
  & Elements overlap or are clearly mis\-aligned; no discernible
    hierarchy; inconsistent spacing.
  & Major regions are aligned and separated; hierarchy is present but
    inconsistent in places.
  & Clear, consistent grid; alignment within and across regions;
    obvious primary/secondary/tertiary hierarchy. \\[2pt]
\textbf{D2} Information Clarity
  & Labels missing or ambiguous; key text un\-readable; no logical
    grouping of controls.
  & Most labels present and legible; grouping is mostly consistent
    with content meaning.
  & All labels descriptive and legible; controls and information
    grouped in a way that makes purpose obvious. \\[2pt]
\textbf{D3} Requirement Fidelity \textit{(Track~B)}
  & Less than 30\% of required elements present; key functionality
    absent.
  & Most required elements present; functionality partially correct
    or with minor omissions.
  & All required elements present; functionality matches the
    requirement without omissions or unrequested additions. \\[2pt]
\textbf{D4} Interaction Quality \textit{(Track~B, dynamic)}
  & Required controls do not respond, or trigger incorrect state
    transitions; task cannot be completed.
  & Task can be completed; some feedback or transition issues
    (e.g.\ missing confirmation, unexpected resets).
  & Task completes smoothly with appropriate feedback at each step;
    state transitions match expectations. \\
\bottomrule
\end{tabular}
\caption{Abbreviated anchor descriptions for the four rubric
dimensions. The full anchor table, including descriptions for scores 2
and 4, appears in Appendix~\ref{appendix:annotation}.}
\label{tab:rubric-anchors}
\end{table}

\subsection{Post-Pilot Adjustment Rule}

A dimension whose pilot Krippendorff $\alpha$ falls below 0.4 is
flagged for redesign before the main annotation begins. The default
remediation is to rewrite ambiguous anchors using examples drawn from
the pilot disagreements. If a second pilot still yields $\alpha < 0.4$,
the dimension is merged with an adjacent one or removed; this
contingency is recorded in Section~\ref{sec:pilot}.

% ----------------------------------------------------------------
\section{Human Annotation Protocol}
\label{sec:annotation-protocol}
% ----------------------------------------------------------------

\subsection{Annotator Recruitment}

The primary annotator pool consists of three peers with prior
coursework in human--computer interaction or visual design, recruited
from the author's department. This pool is sufficient for the planned
sample size and avoids the procurement and quality-control overheads of
crowdsourcing. A contingency plan to recruit additional annotators
through Prolific at \$16/hour, following comparable web-interface
annotation practice in Design2Code~\cite{design2code}, is held in
reserve in case any primary annotator becomes unavailable or the
inter-rater reliability targets are not met after a second pilot.

\subsection{Annotator Instructions}

Annotators receive a written guide (Appendix~\ref{appendix:annotation})
covering: the purpose of the study, the four rubric dimensions and
their full anchor descriptions, three worked examples per dimension,
guidance on tie-handling, and the escalation procedure for ambiguous
items. The guide is reviewed by all annotators in a single 30-minute
session before the pilot, and questions raised in that session are
folded into a frequently-asked-questions appendix to the guide before
the main annotation.

\subsection{Annotation Forms}

\paragraph{Rubric form (Track~B).}
For each item the annotator sees the requirement text, the rendered
screenshot, and an embedded preview of the executable page. The form
collects four ordinal scores (one per dimension), an optional one-line
free-text justification, and a flag for any item the annotator wishes
to escalate. Items flagged by at least two annotators are reviewed by
the author and either retained with a note or discarded.

\paragraph{Pairwise form (subset of Track~B).}
A small subset of Track~B items is also collected as pairwise
comparisons across generator models on the same requirement. This
provides a pairwise reference signal for direct comparison with the
Track~A pairwise judgements, addressing the secondary objective in
RQ2.

\subsection{Tie-Handling and Ambiguity Escalation}

Annotators are instructed to use the centre of the scale (3) only when
the dimension is genuinely intermediate, and to break ties in favour of
the lower score when in doubt. Items receiving more than two escalation
flags across the annotator pool are removed from the reference set and
documented as ambiguous in the benchmark release notes.

% ----------------------------------------------------------------
\section{Inter-Rater Reliability}
\label{sec:irr}
% ----------------------------------------------------------------

Inter-rater reliability is reported per dimension using Krippendorff's
$\alpha$ on the ordinal scale. The targets adopted in this thesis are:

\begin{itemize}
  \item $\alpha \geq 0.7$ — \emph{good}; the dimension is suitable for
    quantitative comparison without further treatment.
  \item $0.5 \leq \alpha < 0.7$ — \emph{acceptable}; results for the
    dimension are reported but accompanied by a caveat about
    annotator variance.
  \item $\alpha < 0.5$ — \emph{insufficient}; the dimension triggers
    the post-pilot adjustment rule of Section~\ref{sec:rubric}.
\end{itemize}

The reference label used in all subsequent analyses is the median of
the three annotator scores per dimension; the mean is reported in
parentheses where it differs by more than half a point.

% ----------------------------------------------------------------
\section{Pilot Annotation}
\label{sec:pilot}
% ----------------------------------------------------------------

\subsection{Pilot Set}

The pilot consists of 10 Track~B items selected to span the range of
expected quality (three high-quality, four mid-quality, three
low-quality, judged by author inspection). All three primary
annotators score every pilot item, yielding $10 \times 3 \times 4 = 120$
individual ratings.

\subsection{Procedure}

Each annotator completes the pilot independently, recording start and
end times per item. The author then computes per-dimension
Krippendorff~$\alpha$ and identifies items on which any two annotators
differ by more than two points; these items become the basis for a
post-pilot calibration discussion.

\subsection{Pilot Outcomes \textit{(to be filled after pilot run)}}

\begin{table}[ht]
\centering
\small
\begin{tabular}{@{}lccc@{}}
\toprule
\textbf{Dimension} & \textbf{Krippendorff $\alpha$} & \textbf{Median time / item} & \textbf{Action} \\
\midrule
D1 Visual Structure       & \textit{TBD} & \textit{TBD} & \textit{TBD} \\
D2 Information Clarity    & \textit{TBD} & \textit{TBD} & \textit{TBD} \\
D3 Requirement Fidelity   & \textit{TBD} & \textit{TBD} & \textit{TBD} \\
D4 Interaction Quality    & \textit{TBD} & \textit{TBD} & \textit{TBD} \\
\bottomrule
\end{tabular}
\caption{Pilot annotation results (placeholder; populated after the
first pilot run completes).}
\label{tab:pilot-irr}
\end{table}

\subsection{Refinements Applied}

This subsection documents any anchor rewrites, dimension merges, or
guide additions made between the pilot and the main annotation, and
re-reports $\alpha$ on a second pilot batch when applicable. The
final, post-pilot rubric is the version used for all reference labels
in Chapters~\ref{chap:experiments}--\ref{chap:results}.
